% =====================================================================
% NegoPlay - Preliminary Report
% Author: Anna Ben-Shushan
% Supervisor: Dr. Yoram Segal
% Course: AI Development Expert
% =====================================================================
\documentclass[11pt, a4paper, oneside]{report}

% ----- Encoding & Fonts -----
\usepackage[T1]{fontenc}
\usepackage[utf8]{inputenc}
\usepackage{lmodern}
\usepackage{microtype}

% ----- Page geometry -----
\usepackage[margin=2.5cm]{geometry}
\usepackage{setspace}
\onehalfspacing

% ----- Tables & Figures -----
\usepackage{graphicx}
\usepackage{booktabs}
\usepackage{array}
\usepackage{tabularx}
\usepackage{longtable}
\usepackage{multirow}
\usepackage{colortbl}
\usepackage[table]{xcolor}
\usepackage{float}

% ----- Math & Symbols -----
\usepackage{amsmath, amssymb}

% ----- Lists -----
\usepackage{enumitem}

% ----- Hyperlinks -----
\usepackage[hidelinks, colorlinks=true, linkcolor=black, citecolor=blue, urlcolor=blue]{hyperref}

% ----- Bibliography -----
\usepackage[round, authoryear]{natbib}
\bibliographystyle{plainnat}

% ----- Header / Footer -----
\usepackage{fancyhdr}
\pagestyle{fancy}
\fancyhf{}
\fancyhead[L]{NegoPlay -- Preliminary Report}
\fancyhead[R]{\leftmark}
\fancyfoot[C]{\thepage}
\renewcommand{\headrulewidth}{0.4pt}

% ----- Custom colors for the WP/Gantt table -----
\definecolor{wpdone}{HTML}{8FCB9B}
\definecolor{wpplan}{HTML}{F8C471}
\definecolor{wpcrit}{HTML}{E74C3C}
\definecolor{wplight}{HTML}{D6EAF8}
\definecolor{headbg}{HTML}{2C3E50}

% ----- Chapter style -----
\usepackage{titlesec}
\titleformat{\chapter}[hang]
  {\normalfont\Huge\bfseries}{\thechapter}{20pt}{\Huge}

% =====================================================================
% DOCUMENT START
% =====================================================================
\begin{document}

% ---------------- Cover Page ----------------
\begin{titlepage}
\centering
\vspace*{1.5cm}

{\Large\bfseries AI Development Expert -- Final Project\par}
\vspace{0.4cm}
{\large Preliminary Report\par}
\vspace{2.0cm}

\rule{\textwidth}{1.5pt}\\[0.6cm]
{\Huge\bfseries NegoPlay\par}
\vspace{0.5cm}
{\Large\itshape Strategic Profiles from Bridge to Business:\\[0.2cm]
An Empirical Investigation of Decision-Making Styles\\[0.2cm] using LLM Agents\par}
\rule{\textwidth}{1.5pt}\\[1.5cm]

\vfill

{\large
\begin{tabular}{rl}
\textbf{Student:}      & Anna Ben-Shushan \\[0.25cm]
\textbf{Supervisor:}   & Dr. Yoram Segal  \\[0.25cm]
\textbf{Course:}       & AI Development Expert \\[0.25cm]
\textbf{Document:}     & Preliminary Project Report \\[0.25cm]
\textbf{Version:}      & 1.0 \\[0.25cm]
\textbf{Date:}         & \today \\
\end{tabular}
}

\vspace{1.5cm}
\end{titlepage}

% ---------------- Dedication / Acknowledgements ----------------
\thispagestyle{empty}
\vspace*{4cm}
\begin{center}
{\Large\bfseries Acknowledgements}
\end{center}
\vspace{1cm}
\noindent
I would like to thank \textbf{Dr.\ Yoram Segal}, my project supervisor, for his structured methodology framework -- in particular the requirement of clearly identified anchor papers, single-SDK architecture, and explicit baseline comparisons. These principles shaped this project from the first day.

\vspace{0.6cm}
\noindent
This project is also the empirical baseline for my PhD research at LUT University (Finland), under the supervision of \textbf{Prof.\ Jari H\"am\"al\"ainen}, and benefits from ongoing discussions with my advisor \textbf{Nezer} on risk-taking behavior in elite bridge play.

\vspace{0.6cm}
\noindent
The 149{,}208-row dataset used in this work was collected from the public EuroBridge database (European Bridge League, 2016--2025) and processed using a custom Python pipeline developed during the first weeks of 2026.

\clearpage

% ---------------- TOC, LoF, LoT ----------------
\tableofcontents
\clearpage
\listoffigures
\listoftables
\clearpage

% =====================================================================
% CHAPTER 1 -- INTRODUCTION
% =====================================================================
\chapter{Introduction}
\label{ch:intro}

\section{Background and Motivation}

Decision making under incomplete information is one of the hardest problems in artificial intelligence. While modern AI has reached super-human performance in fully observable games such as chess \citep{silver2017alphazero}, Go, and Atari \citep{schrittwieser2019muzero}, real-world strategic environments -- business negotiation in particular -- are characterized by hidden information, partial communication, and a mix of cooperation and competition that is hard to model directly.

\textbf{Bridge} is a unique testbed for these challenges. Each of the four players sees only 13 of the 52 cards, must cooperate with a partner while competing against two opponents, and communicates through a highly structured bidding language with approximately $10^{47}$ possible sequences. Bridge bidding is, formally, a multi-agent negotiation under imperfect information with measurable outcomes. \citet{rong2019competitive} showed that a deep learning system splitting bidding into estimation (ENN) and policy (PNN) networks can beat the world-leading rule-based engine Wbridge5; \citet{brown2019pluribus} demonstrated that imperfect-information AI can defeat elite humans in multi-player poker.

\textbf{NegoPlay} asks a more applied question: \emph{can we extract decision-making profiles from a large dataset of elite bridge play, and reuse those profiles to predict and simulate behavior in business negotiation scenarios?}

\section{Why this matters}

Business negotiation suffers from a fundamental data scarcity problem. Real negotiation transcripts are proprietary, outcomes are rarely measurable in standardized form, and no public datasets allow training of behavioral models. Bridge, in contrast, offers a clean laboratory:

\begin{itemize}[leftmargin=*]
    \item More than 149{,}000 structured decision points with measurable outcomes.
    \item Clear win/loss signals (contract success, IMP scores, VP conversion).
    \item Tournament-grade decisions by expert players, with explicit player identity per seat.
\end{itemize}

If profiles discovered in this clean environment transfer measurably to the noisier domain of negotiation, this opens a methodological path for studying business strategy without proprietary negotiation data.

\section{Literature Review and Positioning}
\label{sec:litreview}

This section grounds NegoPlay in the wider research conversation. Following \citet{lockett2007evolving}, we treat an unknown counterpart as a mixture of \emph{cardinal strategies} that can be discovered and re-used. We organize the literature into five themes.

\subsection*{Theme 1: Player-style clustering from gameplay sequences}

Our methodological anchor is \citet{talwadker2022cognitionnet}, who introduced \textbf{CognitionNet}. Using a sequence-to-sequence interpreter coupled to a classifier, they discovered behavioral profiles directly from raw Rummy telemetry. NegoPlay replaces ``Rummy telemetry'' with ``bridge bidding sequences'' but keeps the same logic: \emph{styles emerge from sequences of actions, without predefined rules.} \citet{sztyber2023bridgehand} provide an alternative path -- learning a continuous 8-dimensional embedding of bridge hands -- that we may incorporate in a second iteration.

\subsection*{Theme 2: Inference + decision split (theory of mind)}

Our second anchor is \citet{rong2019competitive}, whose ENN/PNN architecture separates the inference of partner's hidden cards from the action policy. \citet{wang2024inference} extend this with a conditional GAN inference module. NegoPlay applies the same split at the LLM level: the cluster profile is the \emph{inference} (``what kind of agent am I facing?''), and the system prompt is the \emph{policy} (``how do I act?''). This theory-of-mind-before-action structure is now standard in published bridge AI.

\subsection*{Theme 3: Opponent as mixture of cardinal styles}

\citet{lockett2007evolving} is the theoretical foundation for the entire project. Their result -- that explicit mixture coefficients over a small set of cardinal opponent types outperform monolithic networks -- justifies the central NegoPlay premise: every negotiation counterpart can be treated as a weighted mixture of the few profiles discovered in the bridge data.

\subsection*{Theme 4: Bridging structured games and natural-language negotiation}

\citet{james2021attitude} encode both actions and ``cheap talk'' messages into unified attitude vectors, then cluster them. This is the conceptual glue between the bridge stage (structured signals) and the negotiation stage (free-form text), and it is the paper we cite when defending the transfer claim.

\subsection*{Theme 5: Existing bridge AI engines as benchmarks}

\citet{yeh2018automatic} and \citet{kita2024simple} train deep RL bidders that beat Wbridge5; the open-source BEN engine \citep{lorserker_ben} provides a neural-network-based, scriptable opponent that we can use to ground LLM-vs-LLM games against an objective player. \citet{jarosz2022review} provides a clean survey of the field.

\subsection*{What others did, and what NegoPlay does differently}

\begin{table}[H]
\centering\small
\caption{NegoPlay positioning against representative prior work.}
\label{tab:positioning}
\begin{tabularx}{\textwidth}{p{3.2cm} X X}
\toprule
\textbf{Reference} & \textbf{What they did} & \textbf{What NegoPlay does differently} \\
\midrule
\citet{talwadker2022cognitionnet} & Clustered Rummy players from telemetry sequences. & Cluster bridge \emph{bidding} sequences; then \emph{reuse} the profiles in a second, different domain (negotiation). \\
\citet{rong2019competitive} & Built a deep ENN/PNN bidding system that beats Wbridge5. & Replace the deep networks with \emph{LLM agents} prompted by the discovered profile; trade peak strength for explainability and cross-domain reuse. \\
\citet{lockett2007evolving} & Evolved explicit opponent models using NEAT. & Replace the evolved cardinal set with a \emph{data-driven, statistically validated} cardinal set (K-means + HDBSCAN). \\
\citet{james2021attitude} & Clustered attitudes from actions + cheap talk in toy games. & Apply the same logic to a real, large-scale, expert dataset -- and explicitly test transfer between two domains. \\
\bottomrule
\end{tabularx}
\end{table}

% =====================================================================
% CHAPTER 2 -- RESEARCH QUESTION & OBJECTIVES
% =====================================================================
\chapter{Research Question and Objectives}
\label{ch:rq}

\section{Primary Research Question}

\begin{quote}
\itshape Can LLM agents built from automatically-discovered bridge decision-making profiles exhibit behavioral consistency (\emph{Spearman} $\rho \geq 0.7$) between winning in bridge bidding simulations and winning in parallel business negotiation simulations?
\end{quote}

\section{Hypotheses}

\begin{description}[leftmargin=2cm]
    \item[H1 -- Statistical:] K-means clustering on five behavioral features will produce 3--5 stable profiles with Silhouette score $\geq 0.4$ and $p < 0.05$ against random clustering.
    \item[H2 -- Behavioral:] LLM agents instantiated with these profiles will exhibit measurably different behaviors in bridge bidding (per-profile win-rate variance $> 15\%$).
    \item[H3 -- Transfer (main):] The win-rate ranking across profiles will correlate $\rho \geq 0.7$ (Spearman) between the bridge and negotiation domains.
    \item[Null:] No correlation $> 0.5$ between domains. This too is a valid research finding.
\end{description}

\section{Objectives}

The project has four operational objectives:

\begin{enumerate}[leftmargin=*]
    \item \textbf{Statistical:} discover 3--5 stable decision profiles from $\approx$\,1{,}500 elite players.
    \item \textbf{Linguistic:} translate each profile into a 5--7 trait skill description using an LLM.
    \item \textbf{Engineering:} build four LLM agents (one per profile) with verifiable, structured output.
    \item \textbf{Empirical:} run dual simulations (bridge + negotiation) and measure cross-domain alignment.
\end{enumerate}

\section{Anchor papers}

Following Dr.\ Yoram Segal's methodology framework, every project declares its anchor papers explicitly:

\begin{itemize}[leftmargin=*]
    \item \textbf{Anchor 1 (Stage 1 -- clustering):} \citet{talwadker2022cognitionnet}.
    \item \textbf{Anchor 2 (Stage 3 -- inference/decision split):} \citet{rong2019competitive}.
    \item \textbf{Theoretical foundation:} \citet{lockett2007evolving}.
\end{itemize}

% =====================================================================
% CHAPTER 3 -- METHODOLOGY
% =====================================================================
\chapter{Methodology}
\label{ch:method}

\section{Overview}

NegoPlay is a four-stage hybrid pipeline that combines classical machine learning with large language models (see Figure~\ref{fig:pipeline}).

\begin{figure}[H]
\centering
\fbox{\parbox{0.92\textwidth}{
\centering
\vspace{0.4cm}
\textbf{Stage 1} -- Profile Discovery (K-Means + HDBSCAN on player features)\\
$\downarrow$\\
\textbf{Stage 2} -- Skill Extraction (Gemini Flash 2.0 names each profile)\\
$\downarrow$\\
\textbf{Stage 3} -- Agent Construction (4 LLM agents, profile-conditioned prompts)\\
$\downarrow$\\
\textbf{Stage 4} -- Dual Simulation + Cross-Domain Alignment\\
\vspace{0.2cm}
}}
\caption{The four-stage NegoPlay pipeline. Stages 1 and 4 produce the empirical evidence; stages 2 and 3 produce the agents.}
\label{fig:pipeline}
\end{figure}

\section{Dataset}

The dataset was collected by the author from the public EuroBridge database between January and May 2026, using a custom Python scraping pipeline.

\begin{table}[H]
\centering\small
\caption{NegoPlay dataset summary.}
\label{tab:dataset}
\begin{tabular}{lr}
\toprule
\textbf{Metric} & \textbf{Value} \\
\midrule
Total rows (board $\times$ room) & 149{,}208 \\
European Championships covered & 5 (2016--2025) \\
Unique players (by name)        & $\approx$ 2{,}500 \\
Players with $\geq 20$ declared boards & 1{,}572 \\
Player-name coverage (N/S/E/W $\times$ open/closed) & 99.9\,\% \\
Card-holding coverage (52 cards)                    & 54\,\% \\
Bidding-sequence coverage                            & 31\,\% (100\,\% on Herning 2024 + Poznan 2025) \\
\bottomrule
\end{tabular}
\end{table}

\section{Stage 1 -- Profile Discovery (ML)}

Five behavioural features are computed per player from declarations made by that player:

\begin{itemize}[leftmargin=*]
    \item \texttt{slam\_rate} -- percentage of declarations at slam level (6 or 7).
    \item \texttt{success\_rate} -- percentage of contracts made.
    \item \texttt{preempt\_rate} -- percentage of opening bids at level 2 or higher.
    \item \texttt{double\_rate} -- percentage of double calls (penalty + takeout).
    \item \texttt{avg\_risk\_score} -- composite risk index in $[0,10]$.
\end{itemize}

Features are standardized via \texttt{StandardScaler}. K-Means is run for $k \in \{2,3,4,5,6\}$ and the best $k$ is chosen by Silhouette. HDBSCAN is used as an independent validator -- both clusterings must agree to within $\pm 1$ cluster.

\section{Stage 2 -- Skill Extraction (LLM)}

For each cluster, 5--10 chunks of 20--30 representative games are sent to \textbf{Gemini Flash 2.0} with a strict JSON schema. The LLM is instructed to return 5--7 characteristic skills. Skills appearing in $\geq 3$ chunks are kept. The author then assigns a human-readable profile name to each cluster (e.g.\ \emph{Slam Hunter, Insurance Player, Bluffer, Doubler}).

\section{Stage 3 -- Agent Construction}

Each profile becomes an LLM agent through a structured system prompt that contains: identity (profile name), 5--7 traits, domain context (bridge \textit{or} negotiation), output schema (JSON), and 1--2 few-shot examples. Two agent classes are implemented:

\begin{itemize}[leftmargin=*]
    \item \texttt{BridgeAgent.make\_bid(hand, auction)} -- returns a valid bridge bid.
    \item \texttt{NegotiationAgent.respond\_to\_offer(scenario, history, offer)} -- returns a counter-offer plus reasoning.
\end{itemize}

\section{Stage 4 -- Dual Simulation and Alignment}

\textbf{4a Bridge:} 50 randomly generated hands are played by combinations of profile agents; results are scored in IMPs. As a non-LLM benchmark, profile agents can also play one side of a board against the open-source BEN engine \citep{lorserker_ben} sitting on the other side.

\textbf{4b Negotiation:} four scenarios (M\&A, JV equity split, B2B SaaS pricing, government tender) are run for each pair of agents, with up to 10 rounds and an explicit walk-away condition.

\textbf{4c Alignment:} per-profile win rates are computed for both domains; the cross-domain alignment is measured by \emph{Spearman} correlation. A scatter plot is produced; significance is tested via a two-sided permutation test.

\section{Tools and Stack}

\begin{itemize}[leftmargin=*]
    \item \textbf{Language:} Python 3.11 in a virtual environment.
    \item \textbf{IDE:} Google Antigravity (agent-first), with VS Code as fallback.
    \item \textbf{ML libraries:} \texttt{scikit-learn}, \texttt{gensim}, \texttt{shap}.
    \item \textbf{LLM (default):} Google Gemini Flash 2.0 -- cost \$0.30 / M output tokens.
    \item \textbf{LLM (cross-validation):} Anthropic Claude and OpenAI GPT, routed through a single \texttt{llm\_client.py} wrapper. This strengthens the claim that findings are not a single-provider artifact.
    \item \textbf{Source control:} Git + GitHub, public repository.
    \item \textbf{Experiment tracking:} JSON-Lines logs in \texttt{results/llm\_logs/}, with provider, model, tokens, latency, and cost per call.
\end{itemize}

% =====================================================================
% CHAPTER 4 -- SCHEDULE
% =====================================================================
\chapter{Schedule}
\label{ch:schedule}

The project runs for five weeks (ten sessions, two per week). Each week corresponds to a sprint with one major deliverable.

\begin{table}[H]
\centering\small
\caption{Five-week sprint plan.}
\label{tab:sprints}
\begin{tabularx}{\textwidth}{c c X l}
\toprule
\textbf{Sprint} & \textbf{Week} & \textbf{Theme} & \textbf{Deliverable} \\
\midrule
1 & 1 & Research \& Setup        & PRD, literature review, repo skeleton \\
2 & 2 & Stage 1 + Stage 2        & Validated clusters, named profiles \\
3 & 3 & Stage 3                  & Four working LLM agents \\
4 & 4 & Stage 4                  & Bridge + negotiation simulations + alignment \\
5 & 5 & Polish \& Defense        & Final report, presentation, demo video \\
\bottomrule
\end{tabularx}
\end{table}

% =====================================================================
% CHAPTER 5 -- WORK PACKAGES
% =====================================================================
\chapter{Work Packages}
\label{ch:wp}

The project is decomposed into seven work packages (WP). Each WP groups related tasks across one or more sessions.

\section*{WP1 -- Research foundation}
\textbf{Sessions:} 1--2.\quad \textbf{Deliverables:} PRD v1.0, literature review (17 papers + 9 PhD-context), repository structure, two anchor papers identified.\\
\textbf{Tasks:} (T1.1) project scope alignment with supervisor; (T1.2) PRD draft and freeze; (T1.3) literature search and anchor selection; (T1.4) GitHub repository scaffold; (T1.5) Antigravity + Gemini setup; (T1.6) related-work survey of eight bridge-AI repositories.

\section*{WP2 -- Data pipeline (already in place at project start)}
\textbf{Sessions:} pre-project (Jan--May 2026).\quad \textbf{Deliverable:} \texttt{all\_matches\_full.csv} -- 149{,}208 rows.\\
\textbf{Tasks:} EuroBridge HTML scraping; card backfill via BoardAcross; player-name backfill (99.9\,\% coverage). \emph{This WP is closed; the report documents it but no further work is required for NegoPlay.}

\section*{WP3 -- Stage 1: clustering}
\textbf{Sessions:} 3--4.\quad \textbf{Deliverables:} \texttt{player\_features.csv}, \texttt{player\_clusters.csv}, silhouette plot, validation report.\\
\textbf{Tasks:} (T3.1) \texttt{shared/data\_loader.py}; (T3.2) \texttt{stage1\_clustering/features.py}; (T3.3) \texttt{stage1\_clustering/clustering.py} (K-Means + HDBSCAN); (T3.4) \texttt{stage1\_clustering/validation.py} (Silhouette, Davies-Bouldin, permutation $p$-value); (T3.5) exploratory notebook with PCA / t-SNE plots.

\section*{WP4 -- Stage 2: skill extraction}
\textbf{Sessions:} 4.\quad \textbf{Deliverables:} \texttt{skill\_profiles.json}, \texttt{docs/profile\_descriptions.md}.\\
\textbf{Tasks:} (T4.1) \texttt{shared/llm\_client.py} (unified Gemini/Claude/OpenAI wrapper); (T4.2) game chunker; (T4.3) Gemini-based skill extractor with strict JSON schema; (T4.4) cross-chunk aggregator; (T4.5) manual profile naming by the author.

\section*{WP5 -- Stage 3: agents}
\textbf{Sessions:} 5--6.\quad \textbf{Deliverables:} \texttt{BridgeAgent}, \texttt{NegotiationAgent}, full prompt book.\\
\textbf{Tasks:} (T5.1) \texttt{BaseAgent} abstract class; (T5.2) \texttt{BridgeAgent.make\_bid}; (T5.3) \texttt{NegotiationAgent.respond}; (T5.4) prompt library in \texttt{shared/prompts.py}; (T5.5) manual sanity tests confirming behavioral variance $> 30\,\%$.

\section*{WP6 -- Stage 4: simulation and alignment}
\textbf{Sessions:} 7--8.\quad \textbf{Deliverables:} 50 bridge games, 48 negotiation sessions, alignment report (\emph{the key result}).\\
\textbf{Tasks:} (T6.1) \texttt{bridge\_game.py} runner; (T6.2) \texttt{negotiation.py} runner; (T6.3) outcome scoring; (T6.4) optional BEN-engine adapter; (T6.5) \texttt{alignment.py} -- Spearman correlation + permutation test; (T6.6) final scatter plot.

\section*{WP7 -- Polish and defense}
\textbf{Sessions:} 9--10.\quad \textbf{Deliverables:} business-plan section, slide deck, promo video, defense rehearsal.\\
\textbf{Tasks:} (T7.1) business-plan section (TAM/SAM/SOM, token economics); (T7.2) 12--15 slide deck; (T7.3) 60--90 second promo video; (T7.4) README polish; (T7.5) Q\&A document.

% =====================================================================
% CHAPTER 6 -- GANTT
% =====================================================================
\chapter{Gantt Chart and Dependencies}
\label{ch:gantt}

\section{WP $\times$ Session matrix}

Following the structure recommended by Dr.\ Yoram Segal, Table~\ref{tab:wpmatrix} maps every work-package task onto the ten project sessions. Coloured cells indicate the session in which a given task is active. Critical (blocking) tasks are coloured in dark orange; non-blocking tasks in light blue.

\begin{table}[H]
\centering\scriptsize
\caption{Work-package matrix. Each row is one task; columns are the ten project sessions. Dark orange = critical / blocking. Light blue = non-blocking.}
\label{tab:wpmatrix}
\renewcommand{\arraystretch}{1.2}
\begin{tabular}{|p{4.0cm}|c|c|c|c|c|c|c|c|c|c|}
\hline
\rowcolor{headbg!90}\textcolor{white}{\textbf{Task}}
 & \textcolor{white}{\textbf{S1}} & \textcolor{white}{\textbf{S2}} & \textcolor{white}{\textbf{S3}} & \textcolor{white}{\textbf{S4}}
 & \textcolor{white}{\textbf{S5}} & \textcolor{white}{\textbf{S6}} & \textcolor{white}{\textbf{S7}} & \textcolor{white}{\textbf{S8}}
 & \textcolor{white}{\textbf{S9}} & \textcolor{white}{\textbf{S10}} \\
\hline
\multicolumn{11}{|l|}{\textbf{WP1 -- Research foundation}} \\
\hline
T1.1 Scope alignment            & \cellcolor{wpcrit} &                    &                    &                    &                    &                    &                    &                    &                    &                    \\
T1.2 PRD freeze                 & \cellcolor{wpcrit} & \cellcolor{wpcrit} &                    &                    &                    &                    &                    &                    &                    &                    \\
T1.3 Literature review          & \cellcolor{wplight}& \cellcolor{wplight}&                    &                    &                    &                    &                    &                    &                    &                    \\
T1.4 Repo scaffold              & \cellcolor{wpcrit} &                    &                    &                    &                    &                    &                    &                    &                    &                    \\
T1.5 Antigravity + Gemini setup & \cellcolor{wpcrit} &                    &                    &                    &                    &                    &                    &                    &                    &                    \\
T1.6 Related-work survey        & \cellcolor{wplight}& \cellcolor{wplight}&                    &                    &                    &                    &                    &                    &                    &                    \\
\hline
\multicolumn{11}{|l|}{\textbf{WP3 -- Stage 1: clustering}} \\
\hline
T3.1 Data loader                &                    &                    & \cellcolor{wpcrit} &                    &                    &                    &                    &                    &                    &                    \\
T3.2 Player features            &                    &                    & \cellcolor{wpcrit} &                    &                    &                    &                    &                    &                    &                    \\
T3.3 K-Means + HDBSCAN          &                    &                    & \cellcolor{wpcrit} & \cellcolor{wpcrit} &                    &                    &                    &                    &                    &                    \\
T3.4 Validation \& significance &                    &                    &                    & \cellcolor{wpcrit} &                    &                    &                    &                    &                    &                    \\
T3.5 EDA notebook (PCA/t-SNE)   &                    &                    &                    & \cellcolor{wplight}&                    &                    &                    &                    &                    &                    \\
\hline
\multicolumn{11}{|l|}{\textbf{WP4 -- Stage 2: skills}} \\
\hline
T4.1 Unified LLM client         &                    &                    &                    & \cellcolor{wpcrit} &                    &                    &                    &                    &                    &                    \\
T4.2 Game chunker               &                    &                    &                    & \cellcolor{wpcrit} &                    &                    &                    &                    &                    &                    \\
T4.3 Skill extractor (Gemini)   &                    &                    &                    & \cellcolor{wpcrit} &                    &                    &                    &                    &                    &                    \\
T4.4 Aggregator                 &                    &                    &                    & \cellcolor{wpcrit} &                    &                    &                    &                    &                    &                    \\
T4.5 Profile naming             &                    &                    &                    & \cellcolor{wplight}&                    &                    &                    &                    &                    &                    \\
\hline
\multicolumn{11}{|l|}{\textbf{WP5 -- Stage 3: agents}} \\
\hline
T5.1 BaseAgent                  &                    &                    &                    &                    & \cellcolor{wpcrit} &                    &                    &                    &                    &                    \\
T5.2 BridgeAgent                &                    &                    &                    &                    & \cellcolor{wpcrit} & \cellcolor{wpcrit} &                    &                    &                    &                    \\
T5.3 NegotiationAgent           &                    &                    &                    &                    & \cellcolor{wpcrit} & \cellcolor{wpcrit} &                    &                    &                    &                    \\
T5.4 Prompt library             &                    &                    &                    &                    & \cellcolor{wpcrit} & \cellcolor{wpcrit} &                    &                    &                    &                    \\
T5.5 Sanity tests               &                    &                    &                    &                    &                    & \cellcolor{wplight}&                    &                    &                    &                    \\
\hline
\multicolumn{11}{|l|}{\textbf{WP6 -- Stage 4: simulation}} \\
\hline
T6.1 Bridge runner              &                    &                    &                    &                    &                    &                    & \cellcolor{wpcrit} &                    &                    &                    \\
T6.2 Negotiation runner         &                    &                    &                    &                    &                    &                    & \cellcolor{wpcrit} & \cellcolor{wpcrit} &                    &                    \\
T6.3 Outcome scoring            &                    &                    &                    &                    &                    &                    & \cellcolor{wpcrit} &                    &                    &                    \\
T6.4 BEN adapter (optional)     &                    &                    &                    &                    &                    &                    & \cellcolor{wplight}&                    &                    &                    \\
T6.5 Alignment + significance   &                    &                    &                    &                    &                    &                    &                    & \cellcolor{wpcrit} &                    &                    \\
T6.6 Final scatter plot         &                    &                    &                    &                    &                    &                    &                    & \cellcolor{wpcrit} &                    &                    \\
\hline
\multicolumn{11}{|l|}{\textbf{WP7 -- Polish \& defense}} \\
\hline
T7.1 Business-plan section      &                    &                    &                    &                    &                    &                    &                    &                    & \cellcolor{wplight}&                    \\
T7.2 Slide deck                 &                    &                    &                    &                    &                    &                    &                    &                    & \cellcolor{wpcrit} & \cellcolor{wpcrit} \\
T7.3 Promo video                &                    &                    &                    &                    &                    &                    &                    &                    & \cellcolor{wplight}&                    \\
T7.4 README polish              &                    &                    &                    &                    &                    &                    &                    &                    & \cellcolor{wplight}&                    \\
T7.5 Q\&A rehearsal             &                    &                    &                    &                    &                    &                    &                    &                    &                    & \cellcolor{wpcrit} \\
\hline
\end{tabular}
\end{table}

\section{Dependencies and the critical path}

The critical path is sequential: \textbf{WP1 $\rightarrow$ WP3 $\rightarrow$ WP4 $\rightarrow$ WP5 $\rightarrow$ WP6 $\rightarrow$ WP7}. Each downstream WP depends on the output artefact of the previous one:

\begin{itemize}[leftmargin=*]
    \item WP3 cannot start before WP1 freezes the PRD (the feature set must be aligned with the research question).
    \item WP4 needs the clusters from WP3 to chunk games per profile.
    \item WP5 needs the named profiles from WP4 to build prompts.
    \item WP6 needs the agents from WP5 to run simulations.
    \item WP7 needs the alignment result from WP6 to write the conclusions.
\end{itemize}

\noindent
\textbf{Non-blocking (parallelisable) tasks} are coloured in light blue in Table~\ref{tab:wpmatrix}. These include the exploratory notebook (T3.5), the BEN-engine adapter (T6.4), the README polish (T7.4) and the promo video (T7.3). They can slip a session without delaying the critical path.

\section{High-level Gantt}

\begin{table}[H]
\centering\small
\caption{High-level Gantt by sprint.}
\label{tab:hlgantt}
\renewcommand{\arraystretch}{1.4}
\begin{tabular}{|p{3.5cm}|c|c|c|c|c|}
\hline
\rowcolor{headbg!90}\textcolor{white}{\textbf{Work Package}}
 & \textcolor{white}{\textbf{Wk 1}} & \textcolor{white}{\textbf{Wk 2}} & \textcolor{white}{\textbf{Wk 3}} & \textcolor{white}{\textbf{Wk 4}} & \textcolor{white}{\textbf{Wk 5}} \\
\hline
WP1 Research foundation & \cellcolor{wpcrit}  &                    &                    &                    &                    \\
WP3 Clustering          &                     & \cellcolor{wpcrit} &                    &                    &                    \\
WP4 Skills              &                     & \cellcolor{wpcrit} &                    &                    &                    \\
WP5 Agents              &                     &                    & \cellcolor{wpcrit} &                    &                    \\
WP6 Simulation          &                     &                    &                    & \cellcolor{wpcrit} &                    \\
WP7 Polish \& defense   &                     &                    &                    &                    & \cellcolor{wpcrit} \\
\hline
\end{tabular}
\end{table}

% =====================================================================
% CHAPTER 7 -- DELIVERABLES & SUCCESS CRITERIA
% =====================================================================
\chapter{Deliverables and Success Criteria}
\label{ch:deliv}

\section{How will we know we succeeded?}

The success of NegoPlay is defined by three concrete numbers (Table~\ref{tab:kpi}). All three come from the same code base and from runs that are logged with random seeds for full reproducibility.

\begin{table}[H]
\centering\small
\caption{Research key performance indicators.}
\label{tab:kpi}
\begin{tabularx}{\textwidth}{l c X}
\toprule
\textbf{Metric} & \textbf{Target} & \textbf{How measured} \\
\midrule
Silhouette score (Stage 1) & $\geq 0.4$ & Computed on the chosen $k$; visualized in \texttt{results/stage1\_silhouette.png}. \\
Per-profile behavioral variance (Stage 3) & $> 15\,\%$ & Variance of bridge win-rate across the four agents in identical 50-game tournament. \\
Cross-domain Spearman alignment (Stage 4c) & $\geq 0.70$ & Correlation of per-profile win-rate ranks between bridge and negotiation. \\
\bottomrule
\end{tabularx}
\end{table}

\section{Deliverables}

\subsection*{Software}
\begin{itemize}[leftmargin=*]
    \item A public GitHub repository \texttt{NegoPlay/} with a single SDK entry point (\texttt{src/sdk.py}).
    \item Four reusable modules: \texttt{stage1\_clustering}, \texttt{stage2\_skills}, \texttt{stage3\_agents}, \texttt{stage4\_simulate}.
    \item A unified LLM wrapper (\texttt{shared/llm\_client.py}) supporting Gemini, Claude and OpenAI.
\end{itemize}

\subsection*{Data artefacts}
\begin{itemize}[leftmargin=*]
    \item \texttt{data/processed/player\_features.csv} (1{,}572 rows $\times$ 5 features).
    \item \texttt{data/processed/player\_clusters.csv} (with cluster labels and centroids).
    \item \texttt{data/processed/skill\_profiles.json} (4 named profiles, 5--7 skills each).
\end{itemize}

\subsection*{Reports and figures}
\begin{itemize}[leftmargin=*]
    \item \texttt{results/stage1\_silhouette.png} -- the cluster validity plot.
    \item \texttt{results/bridge\_winrates.csv} and \texttt{results/negotiation\_winrates.csv}.
    \item \texttt{results/alignment\_analysis.png} -- the \emph{headline} scatter plot.
    \item \texttt{results/alignment\_report.md} -- the final research answer.
\end{itemize}

\subsection*{Course deliverables}
\begin{itemize}[leftmargin=*]
    \item This preliminary report (PDF, LaTeX source).
    \item Slide deck (12--15 slides).
    \item Promotional video (60--90 seconds).
    \item Business-plan section (TAM/SAM/SOM + token economics).
\end{itemize}

\section{Engineering quality gates}

\begin{itemize}[leftmargin=*]
    \item \textbf{Test coverage:} $\geq 80\,\%$ on Stage 1 (statistics must be reproducible); $\geq 60\,\%$ overall.
    \item \textbf{Cost cap:} hard cap of \$50 across all LLM providers, alert at \$20, target $< \$15$.
    \item \textbf{Reproducibility:} the full pipeline must run end-to-end on a clean clone with one command.
    \item \textbf{Logging:} every LLM call writes provider, model, tokens, latency and cost to \texttt{results/llm\_logs/}.
\end{itemize}

% =====================================================================
% CHAPTER 8 -- RISK MANAGEMENT
% =====================================================================
\chapter{Risk Management}
\label{ch:risks}

\section{Risk register}

\begin{table}[H]
\centering\small
\caption{Top project risks and mitigations.}
\label{tab:risks}
\renewcommand{\arraystretch}{1.25}
\begin{tabularx}{\textwidth}{p{0.5cm} X c c X}
\toprule
\textbf{\#} & \textbf{Risk} & \textbf{Lik.} & \textbf{Imp.} & \textbf{Mitigation} \\
\midrule
R1 & Gemini Flash 2.0 produces inconsistent JSON or weak reasoning for agents. & Med & High & Lower temperature to 0.3; enforce \texttt{response\_mime\_type=application/json}; if still poor, route the affected stage to Claude / GPT through \texttt{llm\_client.py}. \\
R2 & Clusters are not statistically separable (Silhouette $< 0.4$). & Low & High & Pre-screen features in EDA; try alternative feature sets including BridgeHand2Vec embeddings \citep{sztyber2023bridgehand}; fall back to fewer clusters ($k=3$). \\
R3 & Cross-domain alignment falls below the 0.7 target. & Med & Med & Frame any $\rho < 0.5$ as a valid null result; report partial-support discussion if $0.5 \leq \rho < 0.7$; bridge $\rightarrow$ negotiation transfer is still an empirical contribution either way. \\
R4 & LLM agents produce illegal bridge bids or invalid negotiation moves. & Med & High & Add a validator before each action; on rejection, re-prompt with the error in the next turn; cap retries at 3 to avoid infinite loops. \\
R5 & Gemini API rate-limits (15 req/min on free tier) block Stage 4. & Med & Med & Add exponential backoff via \texttt{tenacity}; pre-batch calls; upgrade to paid tier ($\sim$\$5) if needed. \\
R6 & Cumulative LLM cost exceeds the \$50 hard cap. & Low & High & Per-call cost log in \texttt{results/llm\_logs/}; alert at \$20; pause and review prompts before exceeding the cap. \\
R7 & The 5-week timeline is too tight (course + simultaneous PhD work). & High & High & Strict scope discipline: drop LSTM stretch goal first, then promo-video polish; keep critical path intact. \\
R8 & Antigravity IDE bugs (preview, agents) cause work loss. & Low & Med & Commit on every session close; mirror to local VS Code; keep a backup branch tagged at the end of each WP. \\
R9 & Player-name mismatches across competitions distort cluster sizes. & Low & Med & Use fuzzy matching in \texttt{shared/data\_loader.py}; document any merges in \texttt{data/processed/player\_name\_map.csv}. \\
R10 & Negotiation scenarios are perceived as toy / unrealistic. & High & Med & Frame the report explicitly as a ``proof of concept in simulated environment''; cite this as future work for the PhD; ensure scenarios are based on textbook negotiation structures (M\&A, JV, SaaS, tender). \\
\bottomrule
\end{tabularx}
\end{table}

\section{Risk monitoring}

A short \textbf{risk review} is held at the end of every sprint (i.e.\ every two sessions). The risk register is updated in version control, and any new risk that emerges mid-sprint is added immediately. Risks R7 (timeline) and R3 (alignment outcome) are monitored continuously, since both can change project framing late in the schedule.

% =====================================================================
% CHAPTER 9 -- SUMMARY
% =====================================================================
\chapter{Summary}
\label{ch:summary}

\section{What this report covers}

This preliminary report defines NegoPlay -- a five-week project to discover decision-making profiles from 149{,}208 elite bridge tournament hands and to test, via LLM agents, whether the same profiles predict behavior in business negotiation simulations. It anchors the work in two published papers (\citealp{talwadker2022cognitionnet} for clustering; \citealp{rong2019competitive} for the inference/decision split) and a theoretical foundation (\citealp{lockett2007evolving}). It specifies the dataset, the four-stage methodology, the seven work packages, the WP\,$\times$\,session matrix, the critical path, the deliverables, and ten concrete project risks with mitigations.

\section{Expected contribution}

If the cross-domain alignment hypothesis is supported ($\rho \geq 0.7$), NegoPlay will be the first publicly reproducible demonstration that strategic profiles discovered in a clean game-theoretic environment (bridge) measurably predict behavior in a different, less structured domain (business negotiation). This contributes to the long-running PhD research thread \emph{From Bridge Tables to Global Markets} at LUT University.

If the hypothesis is not supported, the project still contributes: (i) a validated set of bridge decision profiles, useful in its own right for bridge analytics and coaching; (ii) a unified, provider-neutral LLM client; and (iii) a clean methodological template for cross-domain behavioral transfer studies.

\section{Beyond the course: link to the PhD}

NegoPlay is the empirical baseline for Paper 1 of the author's PhD thesis. The four discovered profiles, the prompt library and the alignment methodology are designed to be reused in PhD Year 2 with: a larger dataset (BBO Vugraph adds millions of deals), a deeper model (the ENN+PNN+PPO stack described by \citealp{rong2019competitive} and \citealp{kita2024simple}), and real negotiation transcripts (Year 2--3). This continuity between the course project and the PhD is intentional: the same single-SDK code base will carry the work forward.

% =====================================================================
% BIBLIOGRAPHY
% =====================================================================
\bibliography{references}

\end{document}
